\section{Overview}
\label{sec:overview}

This campaign's north star is a general definition and workflow that assigns
to an arbitrary arithmetic scheme one or more \emph{canonical quantum
systems}: kinematics realised as a Hilbert space, observables realised as an
algebra, and, where possible, dynamics realised as automorphisms or as
projective unitary representations of symmetry groups, with fusion
categories as the expected general endpoint. The word \emph{canonical} is a
technical requirement here, not a mood: the assignment must be exhibited as a
functor, every choice on which it depends must be named explicitly, and a
theorem must state exactly what is independent of those choices.

The concrete guiding path fixes a target before the generalities are
attempted: the prime field $\mathbb{F}_p$ of characteristic $p$, its
canonical symplectic space $\mathbb{F}_p \times \mathbb{F}_p$, and the finite
Weyl--Heisenberg system quantizing it. Every general definition this campaign
proposes must reproduce this case exactly, and is tested against it before
anything else. The result is recorded in this lab book, one topic per
section, each statement carrying its full hypotheses, an honest status, and
its provenance (\S\ref{sec:overview-status}); nothing in these pages is
asserted at a strength its evidence does not support.

\subsection*{Status}
\label{sec:overview-status}

The following table records the \textbf{28} rows issued by the campaign's
first increment
(\texttt{briefs/wh-kappa-target.md}): the canonical quantum system attached
to a finite field $\kappa$, with every choice it depends on named. Of these,
\textbf{22} are \statusProved{} (four of them \statusProvedConditional{},
meaning the proved statement is restricted to $p=2$ or to odd $p$ and
displays that restriction in its own text — never abbreviated into the
label), \textbf{4} are \statusSketch{} (an argument exists and is written
down, but has not yet been through this campaign's hostile-critic round),
\textbf{1} is \statusConjecture{}, and \textbf{1} is \statusRefuted{} (with
its surviving weaker statement recorded in the adjacent row). Every
definition D1--D11 of \texttt{definitions.md} and every row below is
restated in full, under a descriptive name, in the sections that follow;
the table records only where.

The current register, including the local-ring results and the separate
field-with-one-element sidequest, has \textbf{135} rows: \textbf{112}
\statusProved{}, \textbf{19} \statusSketch{}, three \statusConjecture{} and
one \statusRefuted{}. Section~\ref{sec:sidequest-f1} compares the
literature on the field with one element, with qubit and tensor-product
examples; seven of its finite-kernel claims are proved and eight supporting
comparisons remain sketches. Section~\ref{sec:f1-operational} develops the
subsystem-composition direction: a positive Hecke family, its symmetric-group
quantum endpoint, and an exact contextual process theorem. Eleven claims
of this operational increment are proved and seven supporting comparisons
remain sketches.

The categorical-limit increment adds thirty-seven proved claims. It gives
continuous and germ categories for the full finite projection completion,
typed quantum instruments and classical history routing, local endpoint
lifts, and finite Born-probability limits. Further sections establish
unitary coboundary exchange, Schur--Weyl trace comparisons, the positive
mirabolic boundary and explicit symplectic composition correspondences.

Sections~\ref{sec:frobenius-hierarchy-algebra} and~\ref{sec:arithmetic-processes}
add the arithmetic Frobenius, code-transfer and higher-gate programme. Its
ten new claims passed capped review. The proposed positive
incidence/phase comparison is recorded separately from these arithmetic fibres.

Section~\ref{sec:frobenius-boundary} adds four proved claims for a positive
conditional Frobenius orbit boundary. It retains cyclic matrix blocks,
independent quantum composition and standard subfield decoder probabilities.
The full coupling to arithmetic multiplication and Fourier remains open;
the four claims passed independent capped review and mechanical repair.

Section~\ref{sec:galois-embedding-registers} gives three proved claims for
finite separable extensions over arbitrary base fields and their Galois
closures. Section~\ref{sec:arithmetic-fourier-charts} develops the uniform
arithmetic Gram sectors, conditional Fourier charts and tower refinements;
its three positive statements passed the capped review. The full mixed limit and
the signed-orbit filtration in Section~\ref{sec:signed-orbit-conjecture}
remain conjectures.

Section~\ref{sec:positive-arithmetic-reference} adds six proved claims for
a common positive arithmetic reference and its finite operational boundary.
It treats every finite-field embedding and full fixed arithmetic protocols,
retaining coherent composition and rare-event information. The actual
characteristic remains named; the interpolation changes reference counting
and does not identify the earlier angle construction with this new theory.

The increment's central finding, proved in \S\ref{sec:polarizing-cocycle},
is that the construction depends on a second datum beyond the field and the
additive character: a \emph{polarizing cocycle}. At odd characteristic this
choice is immaterial and has a canonical representative. At characteristic
two it is not immaterial: the admissible cocycles split into exactly two
$SL_2(\kappa)$-orbits, distinguished by the Arf invariant of an associated
quadratic form, giving two genuinely non-isomorphic Heisenberg groups — a
fact detectable by an operator identity inside the quantum system itself.
Per \texttt{PRD.md}, a sharp result of this shape is a positive structural
statement about the definition, not a negative one.

\begin{longtable}{@{}p{3.3cm}p{6.2cm}p{2.6cm}p{2.6cm}@{}}
\toprule
\textbf{Claim id} & \textbf{Statement (short)} & \textbf{Status} & \textbf{Where} \\
\midrule
\endfirsthead
\toprule
\textbf{Claim id} & \textbf{Statement (short)} & \textbf{Status} & \textbf{Where} \\
\midrule
\endhead
\midrule
\multicolumn{4}{r}{\emph{continued on next page}}\\
\endfoot
\bottomrule
\endlastfoot
\texttt{WH-FORM} & The symplectic form is $\kappa$-bilinear, alternating, nondegenerate; its $\kappa$-linear and $\mathbb{F}_p$-linear isometries coincide, both $SL_2(\kappa)$. & \statusProved & Thm.~\ref{thm:wh-form} \\
\texttt{WH-COMM} & The Weyl commutation relations hold uniformly in $p$, for every admissible cocycle. & \statusProved & Thm.~\ref{thm:wh-comm} \\
\texttt{WH-ALG} & The twisted algebra is simple, dimension $q^2$, centre $\mathbb{C}$, for every cocycle. & \statusProved & Thm.~\ref{thm:wh-alg} \\
\texttt{WH-POL} & Every line is Lagrangian; a Schr\"odinger model is a pair $(L,\chi)$, and all such pairs give isomorphic modules. & \statusProved & Thm.~\ref{thm:wh-pol} \\
\texttt{WH-POL-2} & At $p=2$, characters induced on a line with $Q_\beta\neq0$ are forced to order $4$; no such forcing at odd $p$. & \statusProvedConditional & Thm.~\ref{thm:wh-pol-2} \\
\texttt{WH-SVN} & Uniqueness of the finite Weyl representation: one irreducible unitary representation of dimension $q$, intertwiner unique up to $U(1)$, every automorphism inner. & \statusProved & Thm.~\ref{thm:wh-svn} \\
\texttt{WH-ALG-MAT} & The algebra is isomorphic to $M_q(\mathbb{C})$, by a non-canonical isomorphism. & \statusProved & Thm.~\ref{thm:wh-alg-mat} \\
\texttt{WH-BETA-a} & Admissible polarizing cocycles form a torsor of size $q^3$ under the symmetric bilinear forms. & \statusProved & Thm.~\ref{thm:wh-beta-a} \\
\texttt{WH-BETA-b} & At odd $p$, the cocycle is immaterial: $\omega/2$ is a canonical representative fixing an explicit isomorphism between any two choices. & \statusProvedConditional & Thm.~\ref{thm:wh-beta-b} \\
\texttt{WH-BETA-c} & At $p=2$, no symmetric cocycle exists: the odd-$p$ route to a canonical choice is unavailable. & \statusProvedConditional & Thm.~\ref{thm:wh-beta-c} \\
\texttt{WH-BETA-d} & At $p=2$, $Q_\beta$ is a quadratic form with polar form $\omega$; the resulting element-order profile of $H_\beta(\kappa)$ is computed, and at least two isomorphism classes occur. & \statusProvedConditional & Thm.~\ref{thm:wh-beta-d} \\
\texttt{WH-BETA-TYPE} & \textbf{Heart of the increment.} The Arf invariant classifies $\mathrm{Adm}(\omega)$ into exactly two $SL_2(\kappa)$-orbits at $p=2$, both nonempty, giving exactly two isomorphism classes of Heisenberg group. & \statusSketch & Thm.~\ref{thm:wh-beta-type} \\
\texttt{WH-BETA-EPS} & The Arf type is detected inside the quantum system itself, by an explicit operator identity. & \statusSketch & Thm.~\ref{thm:wh-beta-eps} \\
\texttt{WH-BETA-LEVEL} & The algebra and its Weyl frame do not depend on $\beta$ at any characteristic; the finite level-$\mu_p$ frame does, and only at $p=2$. & \statusSketch & Thm.~\ref{thm:wh-beta-level} \\
\texttt{WH-WEIL-a} & Frame automorphisms surject onto $SL_2(\kappa)$ with kernel $\hat V$, at every characteristic. & \statusProved & Thm.~\ref{thm:wh-weil-a} \\
\texttt{WH-WEIL-b} & At odd $p$, this extension splits explicitly, giving a projective unitary representation of $SL_2(\kappa)$. & \statusProvedConditional & Thm.~\ref{thm:wh-weil-b} \\
\texttt{WH-WEIL-c} & At $p=2$, $\mu_2$-valued phases occur exactly over the orthogonal subgroup $O(Q_\beta)$, proper for the reference cocycle. & \statusProvedConditional & Thm.~\ref{thm:wh-weil-c} \\
\texttt{WH-WEIL-d} & At $p=2$, $\mu_4$ is forced exactly when $O(Q_\beta)$ is proper — always except the single case $q=2$, anisotropic. & \statusProvedConditional & Thm.~\ref{thm:wh-weil-d} \\
\texttt{WH-WEIL-e} & Whether the extension splits at $p=2$ for every $q=2^m$; verified only at $q=2,4$. & \statusConjecture & Conj.~\ref{conj:wh-weil-e} \\
\texttt{WH-SYMM} & The algebra does not remember the $\kappa$-module structure of $V(\kappa)$; its intrinsic $\mathbb{F}_p$-linear symmetry is the strictly larger $Sp_{2m}(\mathbb{F}_p)$ once $m>1$. & \statusProved & Thm.~\ref{thm:wh-symm} \\
\texttt{WH-CHOICE} & Beyond $\kappa$, exactly two data: $\psi$ and $\beta$; the resulting algebras are always isomorphic, but no isomorphism among them is distinguished. & \statusProved & Thm.~\ref{thm:wh-choice} \\
\texttt{WH-CANON} & The projective Hilbert space is canonical in $(\kappa,\psi)$, at fixed $\beta$. & \statusProved & Thm.~\ref{thm:wh-canon} \\
\texttt{WH-CANON-BETA} & The projective Hilbert space is canonical independently of $\beta$ as well. & \statusSketch & Thm.~\ref{thm:wh-canon-beta} \\
\texttt{WH-FUNCT-a} & The trace-normalized character family is not restriction-compatible along field extensions; trivial exactly when $p\mid n$. & \statusProved & Thm.~\ref{thm:wh-funct-a} \\
\texttt{WH-FUNCT-b} & $(\kappa,\psi)\mapsto A_{\psi,\beta_0}(V(\kappa))$ is a functor on character-compatible embeddings. & \statusProved & Thm.~\ref{thm:wh-funct-b} \\
\texttt{WH-FUNCT-b-SEC} & \emph{Refuted:} a Frobenius-compatible character exists for every finite field. & \statusRefuted & Thm.~\ref{thm:wh-funct-b-sec} \\
\texttt{WH-FUNCT-c} & The Frobenius obstruction, exactly: invariance forces $c\in\mathbb{F}_p$, and $\mathbb{F}_{p^p}$ then restricts trivially to the prime field. & \statusProved & Thm.~\ref{thm:wh-funct-c} \\
\texttt{WH-FUNCT-d} & Compatible character systems over an algebraic closure are exactly the restrictions of its characters nontrivial on $\mathbb{F}_p$. & \statusProved & Thm.~\ref{thm:wh-funct-d} \\
\end{longtable}
